% Cover letter - Abcam Associate Data Scientist
\documentclass[letterpaper,11pt]{article}
\usepackage[empty]{fullpage}
\usepackage[hidelinks]{hyperref}
\usepackage[english]{babel}
\usepackage{parskip}
\setlength{\parskip}{6pt}
\usepackage{setspace}
\addtolength{\oddsidemargin}{-0.5in}
\addtolength{\evensidemargin}{-0.5in}
\addtolength{\textwidth}{0.95in}
\addtolength{\topmargin}{-.80in}
\addtolength{\textheight}{1.7in}
\urlstyle{same}
\raggedright
\setstretch{1.0}
\pagestyle{empty}
\begin{document}

Dear Hiring Manager,

I am writing to apply for the Associate Data Scientist position in the Data \& AI team at Abcam. I am completing an MSc in Health Data Science at University College London, where a distinction is expected, after a BSc in Biological Sciences at Imperial College London, and I currently work as a Data Analyst Intern at TikTok Shop UK. This role sits where my degree and my working experience meet: applied machine learning and statistics, on scientif\kern0pt ic data, for people who need to make a decision.

The essential requirements map closely onto what I do. Python is my main language, and I use it daily to explore, clean and transform messy data from several sources before any modelling begins. My MSc dissertation builds cardiovascular risk models on 458,840 UK Biobank participants, comparing logistic regression against three tree ensembles with attention to calibration rather than headline accuracy. I mention calibration deliberately, because a model that ranks well and is badly calibrated is not usable for the decision it was built for.

On the causal inference you list as a plus, that is the core of my dissertation rather than an aside. I used doubly robust causal forests to estimate the ef\kern0pt fect of lifestyle change at an individual level, with overlap weighting to handle positivity, and I tested how far the conclusions could be pushed before they broke using sensitivity analysis including E-values, plus external validation on a second cohort. I delivered the result as a Streamlit application so that people could run scenarios without touching the code.

I have worked with life-science data before this. At Westlake University I analysed multi-omics data on epigenetic change in ageing skin stem cells, and refactored and parallelised the pipeline in Python and R, cutting processing time by around 60 per cent while adding version-controlled scripts so the work could be reproduced. At the University of Science and Technology of China I built a spatial transcriptomics pipeline for medulloblastoma covering spot detection, f\kern0pt iltering and clustering, where tuning the decoding parameters improved decoding accuracy substantially. My undergraduate project ran molecular dynamics simulations of rare genetic variants and analysed hundreds of gigabytes of output through a Python pipeline.

I should set out where I am starting from. My AWS experience is at a basic level, and while I use Git and Docker routinely, I have not owned a production deployment with a mature CI/CD pipeline, so the deployment and monitoring side of the role is where I would need support early on. My statistics, my Python and my ability to explain a technical result to someone who does not want the technical detail are all well established.

I have the right to work in the UK until January 2029 without sponsorship. I am currently based in London and am willing to relocate to Cambridge for an on-site role.

Thank you for considering my application. I would welcome the opportunity to discuss how my background in machine learning, causal inference and life-science data could support the Data \& AI team.

\vspace{2pt}
Yours sincerely,

\vspace{8pt}
Zhuoyan Sun

\end{document}
